\paragraph{About this part.} Language models can now write code, propose features and suggest model
structures. Part~V asks how that ability fits into the research process of the previous parts without
corrupting it. It first separates three different uses of generative models in LOB research
(Section~\ref{sec:three-uses}); then describes \lamd{}, a loop in which a language model proposes
candidate models and a deterministic replay judges them (Section~\ref{sec:lamd}); explains why a
proposal from a language model is not evidence (Section~\ref{sec:not-validation}); distinguishes
searching a fixed space from inventing explanations (Section~\ref{sec:abduction}); and shows how the
number of candidates tried raises the bar any winner must clear (Section~\ref{sec:search}).

% ================================================================
\section{Three uses of generative models, kept apart}
\label{sec:three-uses}
% ================================================================

Generative and language-model technology now enters LOB research in at least three distinct ways,
and they should not be confused.

\begin{table}[!htbp]
\centering\small
\caption{Three uses of generative models in LOB research.}
\label{tab:three-uses}
\begin{tabularx}{\linewidth}{lXXl}
\toprule
Use & What is generated & Examples & In this paper \\
\midrule
LOB generation & synthetic order and message streams & autoregressive message
  models~\citep{nagy2023} and their evaluation benchmarks~\citep{lobbench2025} & adjacent work \\
LOB modelling & predictions and representations of market state & Transformers, state-space models,
  pretrained encoders~\citep{linna2025lobert} & Part~II \\
Model generation & candidate models, signals and hypotheses & \lamd{} & this part \\
\bottomrule
\end{tabularx}
\end{table}

LOB generation produces market trajectories; the generated messages can be processed by a book
simulator to produce synthetic books. Its value lies in simulation, stress testing and data
augmentation, and its difficulty lies in preserving the joint structure of orders, queues,
executions and prices rather than only their marginal statistics. It is a legitimate research area
but not part of the methodology proposed here: consistent with the first principle of
Section~\ref{sec:intro}, candidate models are evaluated only on observed market data.

A fourth use is now large enough to have its own surveys: the language model as a \emph{trading agent}.
\citet{ding2024llmtrading} review LLM agents for financial trading---their architectures, data inputs,
backtest performance and open challenges---and \citet{dong2025llmfinance} survey LLM agents across five
financial domains (data analysis, investment research, trading, investment management and risk
management), cataloguing over 30 financial benchmarks and 20 representative models and the obstacles to
deployment. \citet{hua2026agentic} review \emph{agentic} quantitative trading across five stages---factor
mining, signal discovery, portfolio construction, order execution and risk management---and find that
current systems concentrate on signal discovery, that integration with execution and risk control is
still uncommon, and that strong model or forecasting capability ``does not reliably translate into
trading performance under live market conditions''. That literature asks how a language model can act
in markets. This part asks a narrower
question: how a language model can take part in quantitative hypothesis generation without becoming
either the source of the market data or the source of the evidence that validates a hypothesis.

LOB modelling is the subject of Part~II. A pretrained encoder that can also predict the next
message, such as LOBERT, is used there as a model of the market, not as a source of data.

Model generation is different in kind. The language model does not produce market data or market
forecasts. It produces \emph{hypotheses about how to trade}, which are then evaluated against the
recorded market.

% ================================================================
\section{The \lamd{} loop}
\label{sec:lamd}
% ================================================================

LLM-Assisted Model Discovery (\lamd)\footnote{An earlier description of the same loop used the
name LLM-Assisted Trading Strategy Search (LATSS)~\citep{mayeski_latss}.} is a research loop in which a language model proposes
candidate models and every other step is external to it and reproducible: given the same data,
configuration and random seeds, it gives the same result. It belongs to a broader
family of methods in which a language model is paired with an automated evaluator to search a
space of programs~\citep{romera2024}. The closest published application in finance is
\citet{kou2025strategy}, in which prompted language models generate executable alpha-factor
candidates, agents filter them by market state and predictive quality, and the surviving factors are
combined with dynamically optimised weights; the authors report a 53.17\% cumulative return on the SSE~50
from January 2023 to January 2024, ahead of their benchmarks. The difference from
\lamd{} lies less in the generation step than in what surrounds it: \lamd{} makes the number of
candidates tried part of the evidence (Section~\ref{sec:search}) and, as proposed here for order-book
strategies, evaluates each candidate's execution in queue-exact replay on recorded order-level data; the
existing implementation~\citep{mayeski_latss} works on daily data. The components are:

\begin{enumerate}[leftmargin=1.5em]
\item \textbf{Vocabulary.} A human-specified set of primitives the proposals may use: OFI at several
  windows, queue imbalance, trade flow, Hawkes intensities, returns and volatility, queue position,
  regime indicators, cross-asset inputs, and execution actions (place, keep, cancel, cross).
\item \textbf{Proposal.} The language model proposes or mutates a candidate: a feature combination,
  a transformation, a decision rule or a model structure, expressed in the vocabulary.
\item \textbf{Compilation.} A deterministic compiler turns the proposal into an executable model
  with explicit free parameters. Proposals that do not compile are rejected without evaluation.
\item \textbf{Fitting.} Parameters are fitted on a training window only.
\item \textbf{Replay.} The fitted model drives an execution policy in a queue-exact replay of
  recorded data (Part~IV), against the model-free benchmark of Section~\ref{sec:shadow}.
\item \textbf{Validation.} Walk-forward out-of-sample evaluation with purging and embargo
  (Section~\ref{sec:validation}).
\item \textbf{Certification.} The out-of-sample result is compared against a null that accounts for
  every candidate evaluated so far (Section~\ref{sec:search}).
\item \textbf{Archive.} Every candidate---prompt, compiled model, fit, split, result---is recorded,
  whether it passed or failed, and the archive seeds the next round of proposals. The running count
  of candidates evaluated is $N_{\text{trials}}$.
\end{enumerate}

\begin{figure}[ht]
\centering
\begin{tikzpicture}[
  node distance=3.5mm,
  box/.style={draw, rounded corners=2pt, minimum width=62mm, minimum height=6.5mm, align=center, font=\small},
  llm/.style={box, fill=blue!8},
  det/.style={box, fill=gray!10},
  arr/.style={-{Stealth[length=2mm]}, thick}]
\node[box] (v) {vocabulary (human-specified)};
\node[llm, below=of v] (p) {language model: propose / mutate};
\node[det, below=of p] (c) {deterministic compiler};
\node[det, below=of c] (f) {parameter fit (training window)};
\node[det, below=of f] (r) {queue-exact replay on recorded data};
\node[det, below=of r] (o) {walk-forward OOS, purge, embargo};
\node[det, below=of o] (s) {certification against search-aware null};
\node[det, below=of s] (a) {archive (all candidates, $N_{\text{trials}}$)};
\foreach \x/\y in {v/p,p/c,c/f,f/r,r/o,o/s,s/a} \draw[arr] (\x) -- (\y);
\draw[arr] (a.east) -- ++(8mm,0) |- node[pos=0.25, right, font=\scriptsize]{mutate / iterate} (p.east);
\end{tikzpicture}
\caption{The \lamd{} loop. Only the shaded blue step involves the language model; every step that
touches data or produces evidence is a fixed, reproducible procedure.}
\label{fig:lamd}
\end{figure}

\begin{example}[One hypothetical iteration, three candidates]
In the table, $\OFI^{(5)}_t$ is $\OFI^{(w)}_t$ with $w = 5$ s; $\lambda_{\text{sell}}(t)$ is the
component $\lambda_i(t)$ of a fitted Hawkes model for which $i$ is the sell-market-order type, and
$\bar\lambda_{\text{sell}}$ the corresponding component of $\bar\lambda$; $\xi$ and $c$ are
thresholds fitted on the training window.
\begin{center}\small
\begin{tabularx}{\linewidth}{lXlll}
\toprule
id & proposal (paraphrased) & compiled form & OOS net, ticks/fill & $N_{\text{trials}}$ \\
\midrule
c1 & keep a passive bid only while short-horizon OFI agrees & cancel if $\OFI^{(5)}_t < \xi$ & $+0.04$ & 1 \\
c2 & interact OFI with queue imbalance & logit on $\OFI^{(5)}_t \times \QI_t$ & $+0.01$ & 2 \\
c3 & cancel faster when sell-order intensity spikes & cancel if $\lambda_{\text{sell}}(t) > c\,\bar\lambda_{\text{sell}}$ & $-0.02$ & 3 \\
\bottomrule
\end{tabularx}
\end{center}
Each result is measured relative to the model-free benchmark in the same replay. Candidate c1 is
not ``discovered'' when it tops this table; it is discovered only if it clears a null adjusted to
the total number of candidates the loop has evaluated, which after a few days of automated search
may be in the tens of thousands.
\end{example}

\subsection{A worked implementation}

The loop has been implemented for daily-frequency strategies in the LATSS system described
in~\citet{mayeski_latss}, and the details show which choices matter.
\begin{itemize}[leftmargin=1.5em,itemsep=1pt]
\item \textbf{The language model writes structure, never sees prices.} It composes transformations from
  a fixed set of fifteen tools (moving averages, rate of change, z-scores, realised volatility,
  breakouts, percentile ranks, crossovers and similar). Prices reach the code only when the engine
  runs it. Everything the system can discover is some arrangement of those tools.
\item \textbf{The model declares parameters; an optimiser sets them.} Every window and threshold is a
  declared ``knob'' with a range; differential evolution chooses its value on training data.
\item \textbf{Validation is fold-based.} Data are split by quarter; candidates are fitted on some
  quarters and scored on the held-out rest, and the median is kept. The author notes two weaknesses:
  long look-back windows leak a little information across fold boundaries, and random quarters
  resample a single regime, which missed a later regime change.
\item \textbf{The skill bar is empirical.} Each candidate's positions are circularly shifted in time,
  which keeps their exposure and turnover but scrambles their timing; the best score achieved by
  these skill-free shifted versions is the bar a candidate must clear.
\end{itemize}
The author's own assessment is the right one for this survey: with a generic vocabulary the results
are generic, and the human decisions---what to look for, which tools to supply, how hard to test, what
to trust---remain the binding constraints. For LOB models the vocabulary would be the primitives of
Part~II, and the evaluator the queue-exact replay of Part~IV.

% ================================================================
\section{Hypothesis generation is not validation}
\label{sec:not-validation}
% ================================================================

A candidate proposed by a language model has no stronger evidentiary status than one proposed by a
person. Its plausibility, its fluency, and the quality of the explanation that accompanies it are
not evidence that it captures market structure. The scientific content of \lamd{} is in the
deterministic steps around the proposal, not in the proposal.

The same caution applies to language models used directly on market data.
\citet{llmlob2026} train a language model on synthetic LOB event sequences and find that it
generates valid sequences almost perfectly, yet its implicit world model fails to learn the state
of the book, which leads to biased estimates and spurious predictability when it is used to
forecast events. Syntactic validity is not evidence of a learned mechanism, and the same holds for
plausible-looking trading rules. Only external evaluation on held-out observed data can establish that a
candidate works.

% ================================================================
\section{Search versus abduction}
\label{sec:abduction}
% ================================================================

A distinction from the philosophy of science clarifies what such a system can and cannot do.
\emph{Search} explores a given space of hypotheses: ``here is a space of possible solutions; try
things''. \emph{Abduction} creates a hypothesis that was not in the space: ``something about reality
does not make sense; what new concept would make sense of it?''~\citep{mayeski2026abduction}. Language
models are effective at search and recombination: program-search systems built on them have found new
mathematical constructions~\citep{romera2024}, and \lamd{} finds new arrangements of supplied tools.
\citet{zahavy2026} argues that current language models handle induction and, increasingly,
deduction, but lack a mechanism for the abductive step from experience to new principles.

For order-book research the implication is concrete. A \lamd{}-style system can search exhaustively
over combinations of known primitives---OFI at several windows, queue imbalance, Hawkes intensities,
queue position, state modulators, targets and horizons---far faster than a person. A system restricted
to such a vocabulary cannot discover a primitive outside it: a new state variable, an unrecognised
interaction between hidden and visible liquidity, or a new relation between depletion and adverse
selection. A less constrained system could, in principle, search over representations as well as over
combinations of supplied primitives; whether current language models can do so reliably is an empirical
question. Whether that limitation is permanent is an open question; that it holds for present systems
is the working assumption of this survey~\citep{mayeski2026agent}. Automated search is an amplifier of
human-specified structure, and its outputs must pass the same certification as any human proposal.

% ================================================================
\section{Search complexity and multiple testing}
\label{sec:search}
% ================================================================

The more candidates a search evaluates, the larger the best in-sample result it will find by
chance. For $N_{\text{trials}}$ independent candidates with no skill, the expected maximum of
their standardised Sharpe-ratio estimates grows roughly like $\sqrt{2 \ln N_{\text{trials}}}$. (The Sharpe ratio is
average return divided by the standard deviation of return---profit per unit of risk.) In words:
toss a thousand fair coins a hundred times each, and the luckiest coin will look like a skilled
one; the more coins, the luckier the best one looks, though the growth is slow. The deflated Sharpe
ratio~\citep{bailey2014} makes the correction precise and accounts for non-normal returns
(Section~\ref{sec:generalisation}). Related tools include the reality check for data
snooping~\citep{white2000}, the probability of backtest overfitting~\citep{bailey2014pbo}, and
raised significance thresholds for factor discovery~\citep{harvey2016}.

\begin{example}[The same Sharpe, different evidence]
Take the standard error of an annualised Sharpe estimate as approximately $1/\sqrt{T_{\text{yr}}}$ for $T_{\text{yr}}$
years of data, so the null maximum is about $\sqrt{2\ln N_{\text{trials}}}/\sqrt{T_{\text{yr}}}$:
\begin{center}\small
\begin{tabular}{rcc}
\toprule
$N_{\text{trials}}$ & null maximum, 1 year & null maximum, 4 years \\
\midrule
1 & 0 & 0 \\
100 & $\approx 3.0$ & $\approx 1.5$ \\
50\,000 & $\approx 4.65$ & $\approx 2.3$ \\
\bottomrule
\end{tabular}
\end{center}
A Sharpe ratio of 2 from a single pre-specified test on four years is evidence. The same Sharpe
ratio as the best of 50\,000 automatically generated variants on four years is below what noise
alone is expected to produce.
\end{example}

Automated proposal therefore changes the evaluation problem from ``performance and model
complexity'' to
\begin{equation}
  \boxed{\ \text{performance} \;+\; \text{model complexity} \;+\; \text{search complexity}\ }
\end{equation}
\inwords{to judge a discovered model you need three things, not one: how well it did, how
complicated it is, and how many other candidates were tried before it was found.}
The archive's trial count is as much a part of the result as the out-of-sample Sharpe ratio.
Faster replay makes this worse, not better: an order-of-magnitude increase in simulation speed is an
order-of-magnitude increase in the number of hypotheses that can be tested, and the skill threshold
must rise with it.

The central research question for automated model discovery is thus not whether an AI system can
generate a profitable-looking model, but whether an automated search can discover economically
significant structure while controlling the false-discovery rate that the search itself induces.
